<%= heading2 %>{Handling of paragraphs}

<%- bexample %>
One sentence per line.
This rule is important for the usage of version control systems.
A new line is generated with a blank line.
As you would do in Word:
New paragraphs are generated by pressing enter.
In LaTeX, this does not lead to a new paragraph as LaTeX joins subsequent lines.
To start a new paragraph, press enter twice.
This leads to an empty line.
In word, there is the functionality to press shift and enter.
This leads to a hard line break.
The text starts at the beginning of a new line.
In LaTeX, you can do that by using two backslashes (\textbackslash\textbackslash).
\\
Running text seldom needs a hard line break.

Please do \textit{not} use two backslashes for new paragraphs.
For instance, this sentence belongs to the same paragraph, whereas the last one started a new one.
Andrew Stacey gives a long motivation for that in the section <%- bquote %>Line Wrapping<%- equote %> of his \href{https://loopspace.mathforge.org/HowDidIDoThat/TeX/VCS/#section.3}{\emph{Document Revision System}}\footnote{\url{https://loopspace.mathforge.org/HowDidIDoThat/TeX/VCS/}}.
<%- eexample %>

Once the text is final, single paragraphs can be polished.
If the last line of a paragraph holds only a word or two, or is nearly full, let TeX set the paragraph one line shorter or longer: write \texttt{\textbackslash looseness=-1} (or \texttt{=1}) as the first thing in that paragraph.
TeX resets \texttt{\textbackslash looseness} at the end of every paragraph, so only that one paragraph is affected.
With LuaLaTeX, a lone short word as the last line is avoided automatically (package impnattypo).
